% A first draft for Computational Statistics. The journal uses the Springer Nature
% template (sn-jnl); the article class is kept here so that the draft compiles anywhere.
\documentclass[11pt,a4paper]{article}
\usepackage[utf8]{inputenc}
\usepackage[T1]{fontenc}
\usepackage{lmodern}
\usepackage{amsmath,amssymb,amsthm}
\usepackage{booktabs}
\usepackage[round]{natbib}
\usepackage[margin=2.5cm]{geometry}
\usepackage[colorlinks=true,linkcolor=blue,citecolor=blue,urlcolor=blue]{hyperref}

\newcommand{\vect}[1]{\mathbf{#1}}
\newcommand{\E}{\mathrm{E}}
\DeclareMathOperator{\KL}{KL}
\DeclareMathOperator{\Var}{Var}
\DeclareMathOperator{\Cov}{Cov}

\newtheorem{proposition}{Proposition}
\newtheorem{corollary}{Corollary}
\newtheorem{lemma}{Lemma}
\theoremstyle{remark}
\newtheorem{remark}{Remark}
\newtheorem{example}{Example}

\sloppy

\title{Estimating exponential smoothing with unobserved periods: the observed-data
likelihood of ADAM and iETS}
\author{Ivan Svetunkov\thanks{TODO: affiliation and co-authors.}}
\date{Draft, October 2026}

\begin{document}
\maketitle

\begin{abstract}
Exponential smoothing in its single source of error (SSOE) form is estimated by running the
model's recursion through the data and maximising the likelihood of the one-step-ahead errors.
This is the exact likelihood when every observation is available, because the states are then
known functions of the data and the parameters. We show that it stops being so as soon as some
periods are unobserved: the zeros of an intermittent demand model (iETS) and the missing values
in the Augmented Dynamic Adaptive Model (ADAM). The usual remedies do not restore it. The
differential entropy of the unobserved sizes that iETS adds to the likelihood does not follow
from the EM algorithm it is derived from: its stationary points differ from those of the
observed-data likelihood, and it shrinks the estimate of the scale by the share of the observed
periods. Propagating the states through the unobserved periods with zero errors and keeping the
one-step density ignores that the error at the next observation is a multistep one. For the linear
Gaussian models and for the pure multiplicative ETS of ADAM with the log-normal distribution, which
is linear in logarithms, we derive the exact observed-data likelihood: a Kalman filter in which the
state covariance builds up over the unobserved periods and the gain varies with it. It reduces to
the usual SSOE recursion when all periods are observed. For the other pure multiplicative models
we give an approximation that mixes the one-step density over the drift of the states, with its
moments by quadrature, and moment equations for the scale solved by fixed-point iteration. We also
correct the estimator of the scale of the log-normal distribution used in ADAM, which is not the
maximum likelihood one, and the prediction intervals of the models with an occurrence part.
\end{abstract}

\noindent\textbf{Keywords:} exponential smoothing; state space models; intermittent demand;
missing values; Kalman filter; EM algorithm; maximum likelihood.

% =============================================================================
\section{Introduction}
\label{sec:intro}

Exponential smoothing (ETS) in the innovations, or single source of error (SSOE), state space form
\citep{Snyder1985,OrdKoehlerSnyder1997,Hyndman2008} is estimated by a remarkably simple procedure:
the model is run through the data with its own recursion, the one-step-ahead errors are recorded,
and their likelihood is maximised. The simplicity rests on one property. With a single source of
error, the state vector at time $t$ is a deterministic function of the initial states, the
parameters and the observations up to $t$, so that the conditional distribution of each observation
given the past is the distribution of the one-step error. No filter is needed, and the likelihood is
the product of the one-step densities.

The Augmented Dynamic Adaptive Model \citep[ADAM,][]{ADAM} extends this framework to ARIMA,
regressors, multiple seasonalities, a variety of distributions and, through an occurrence part, to
intermittent demand \citep[iETS,][]{iETS}. Two situations in ADAM involve periods in which the
variable of interest is not observed. In intermittent demand the demand is $y_t = o_t z_t$, the
product of a Bernoulli occurrence $o_t$ and a potential demand size $z_t$, and the size is not
observed when $o_t = 0$. With missing values, nothing is observed at those periods. In both cases
the implementation follows the recursion of the complete-data case: the error is set to zero at the
unobserved periods, so that the states go on with the transition alone, and the likelihood is that of
the one-step errors at the observed periods. For the occurrence model, \citet{iETS} add the
differential entropy of the unobserved sizes to the likelihood, as the expected log-density of the
missing values, by analogy with the EM algorithm \citep{EM}.

This paper shows that neither of these gives the likelihood of the observed data, and derives what
does. The argument is structural rather than empirical: we compare the objective functions, their
stationary points and the distributions of the errors, and we do not present a forecasting
evaluation here. The contributions are as follows.
\begin{enumerate}
\item The entropy-augmented likelihood of iETS is not the objective of the EM algorithm: replacing
the cross-entropy of the E-step by the entropy at the current parameters changes the stationary
points. The EM algorithm converges to the stationary points of the observed-data likelihood, in which
the unobserved sizes do not appear at all; the entropy version shrinks the estimate of the scale by
the share of the observed periods and is unbounded with a dynamic scale
(Section~\ref{sec:entropy}).
\item With unobserved periods the states are no longer known functions of the data, and the SSOE
recursion with a fixed gain is not the conditional mean of the states. For linear Gaussian models the
exact observed-data likelihood is given by a Kalman filter written for the SSOE form, whose state
covariance is zero as long as every period is observed, builds up over the unobserved periods and
decays afterwards, and whose gain varies with it (Section~\ref{sec:linear}). The familiar SSOE
recursion is its special case for complete data.
\item The pure multiplicative ETS of ADAM, whose states are updated by powers of $1+\epsilon_t$, is
linear in logarithms, and with the log-normal distribution the same filter applied to the logarithms
gives its exact likelihood, with the scale in closed form. The estimator of the scale of the
log-normal distribution used in ADAM is not the maximum likelihood one, even without unobserved
periods (Section~\ref{sec:multiplicative}).
\item For the other pure multiplicative models (the Gamma and the inverse Gaussian distributions,
and the conventional multiplicative ETS) the logarithms of the states are linear in non-Gaussian
innovations. We approximate the predictive density by the one-step density mixed over a Gaussian
drift of the states, with moments by Gauss--Legendre quadrature on the quantile function and the
mixture by Gauss--Hermite quadrature, and estimate the scale by moment equations that account for
the drift, solved by fixed-point iteration (Section~\ref{sec:approximate}).
\item The multistep structure carries over to the prediction intervals: the interval $h$ steps ahead
is the density after an unobserved stretch of $h$ periods, and the intervals of a model with an
occurrence part are intervals of the sizes, centred on their forecast rather than on that of the
demand (Section~\ref{sec:intervals}).
\end{enumerate}
The computational aspects, and the parts of this programme that are implemented in the
\texttt{smooth} package for R and Python, are discussed in Section~\ref{sec:computation}, and
Section~\ref{sec:conclusions} concludes.

% =============================================================================
\section{The model and its current estimation}
\label{sec:model}

\subsection{ADAM in the SSOE form}

ADAM \citep{ADAM} writes a model with the state vector $\vect v_t$ and the vector of lags
$\vect l$ as
\begin{equation}
z_t = w(\vect v_{t-\vect l}) + r(\vect v_{t-\vect l})\,\epsilon_t, \qquad
\vect v_t = f(\vect v_{t-\vect l}) + g(\vect v_{t-\vect l})\,\epsilon_t ,
\label{eq:adam}
\end{equation}
where $\vect v_{t-\vect l}$ holds each component at its own lag and $\epsilon_t$ are independent
errors. Two classes of models matter here.

\paragraph{Pure additive models.} With the measurement vector $\vect w$, the transition matrix
$\vect F$ and the persistence vector $\vect g$,
\begin{equation}
z_t = \vect w'\vect v_{t-\vect l} + \epsilon_t, \qquad
\vect v_t = \vect F\vect v_{t-\vect l} + \vect g\epsilon_t .
\label{eq:additive}
\end{equation}
This covers ETS with additive components, ARIMA and regressors in ADAM. Stacking the lagged
components gives the usual form with a single lag, $\vect x_t = \vect F_x \vect x_{t-1} +
\vect g_x\epsilon_t$, $z_t = \vect w_x'\vect x_{t-1} + \epsilon_t$ \citep{Hyndman2008};
we use it in the derivations and write $\vect w$, $\vect F$ and $\vect g$ for its matrices.

\paragraph{Pure multiplicative models.} With all the components multiplicative (or absent), the
location is the product of the components, $\mu_t = l_{t-1} b_{t-1}^{\phi} s_{t-m}$, and
$z_t = \mu_t(1+\epsilon_t)$ with $\E(1+\epsilon_t) = 1$. The conventional form of ETS
\citep{Hyndman2008} updates each component $k$ by the factor $1+g_k\epsilon_t$, for instance
$l_t = l_{t-1}b_{t-1}^{\phi}(1+\alpha\epsilon_t)$ and $b_t = b_{t-1}^{\phi}(1+\beta\epsilon_t)$. The
form of ADAM \citep{ADAM} updates it by the power $(1+\epsilon_t)^{g_k}$, for instance
$l_t = l_{t-1}b_{t-1}^{\phi}(1+\epsilon_t)^{\alpha}$.

\subsection{Unobserved periods}

Let $O_t = 1$ when $z_t$ is observed and $O_t = 0$ otherwise. In iETS \citep{iETS},
$y_t = o_t z_t$ with $o_t \sim \text{Bernoulli}(p_t)$ independent of $z_t$, and $O_t = o_t$; with
missing values, $O_t = 0$ at the missing periods. $T$ is the number of periods, $T_1 = \sum_t O_t$
the number of the observed ones, and $j_t$, defined at the observed periods, is the number of
periods since the previous observed one (the demand interval in iETS), with $j_t = 1$ when $z_{t-1}$
is observed and the first interval counted from the initial states.

\subsection{The current estimator}

The implementation of ADAM, and the estimator of \citet[Appendix~C]{iETS}, run the recursion of
\eqref{eq:adam} with $\epsilon_t = 0$ at the unobserved periods:
\begin{equation}
\hat{\vect v}_t = f(\hat{\vect v}_{t-\vect l}) \quad (O_t = 0), \qquad
\hat{\vect v}_t = f(\hat{\vect v}_{t-\vect l}) + g(\hat{\vect v}_{t-\vect l})\,\hat\epsilon_t
\quad (O_t = 1),
\label{eq:recursion}
\end{equation}
with $\hat\epsilon_t$ the error of $z_t$ against $\hat\mu_t = w(\hat{\vect v}_{t-\vect l})$, which is
the $j_t$-steps-ahead point forecast from the states at the previous observed period. The
log-likelihood of the sizes is then $\sum_{O_t=1}\log f_z(z_t\mid\hat\mu_t,\sigma)$ with the
one-step density $f_z$. For the occurrence model, \citet{iETS} add, for each zero, the expected
log-density of the unobserved size, $\int f_z\log f_z\,dz = -H_z$, the negative differential entropy:
\begin{equation}
\ell_H(\theta) = \sum_{O_t=1}\log f_z(z_t\mid\hat\mu_t,\sigma) - \sum_{O_t=0}H_z(\hat\mu_t,\sigma)
+ \sum_{O_t=1}\log p_t + \sum_{O_t=0}\log(1-p_t),
\label{eq:entropy}
\end{equation}
with $\theta$ the parameters of the sizes \citep[eq.~B.6]{iETS}.

% =============================================================================
\section{The observed-data likelihood and the entropy}
\label{sec:entropy}

\begin{proposition}[Observed-data likelihood]
\label{prop:observed}
Let $Y$ be the observed data, $\theta$ the parameters of the sizes and $\lambda$ those of the
occurrence. Then
\begin{equation}
\ell(\theta,\lambda\mid Y) = \sum_{O_t=1}\log f\!\left(z_t\mid Y_{t-1};\theta\right)
+ \sum_{O_t=1}\log p_t + \sum_{O_t=0}\log(1-p_t),
\label{eq:observed}
\end{equation}
where $f(z_t\mid Y_{t-1};\theta)$ is the predictive density of the size given the observed past.
The unobserved sizes do not appear, and the two parts separate.
\end{proposition}

\begin{proof}
Since $z_t$ is continuous and independent of $o_t$, $P(y_t = 0\mid Y_{t-1}) =
(1-p_t)\int f(z\mid Y_{t-1})\,dz = 1-p_t$, and $y_t > 0$ has the density $p_t f(y_t\mid Y_{t-1})$.
The prediction-error decomposition of the joint density of the observed data gives
\eqref{eq:observed}. Equivalently, the complete-data likelihood $L_c(\theta,\lambda\mid Y,Z_0)$ with
the unobserved sizes $Z_0$ \citep[eq.~B.3]{iETS} integrates over $Z_0$ to \eqref{eq:observed}.
\end{proof}

The predictive density in \eqref{eq:observed} is the one-step density only when the states at
$t-1$ are known given $Y_{t-1}$, which is the subject of Section~\ref{sec:linear}. Here we compare
\eqref{eq:entropy} with the EM algorithm that motivates it.

\begin{proposition}[The entropy is not the E-step]
\label{prop:em}
The E-step of the EM algorithm for the complete-data likelihood takes, at each zero, the
cross-entropy $\int f_z(z\mid\theta^{(k)})\log f_z(z\mid\theta)\,dz$ at the current estimate
$\theta^{(k)}$, and its gradient in $\theta$ at $\theta = \theta^{(k)}$ is zero. The fixed points of
EM are therefore the stationary points of \eqref{eq:observed}. The entropy term
$-H_z(\theta) = \int f_z(z\mid\theta)\log f_z(z\mid\theta)\,dz$ of \eqref{eq:entropy} has a non-zero
gradient, and the stationary points of $\ell_H$ differ from those of \eqref{eq:observed}.
\end{proposition}

\begin{proof}
The E-step is $Q(\theta\mid\theta^{(k)}) = \E_{Z_0\mid Y,\theta^{(k)}}\log L_c(\theta\mid Y,Z_0)$,
in which the term of a zero is the cross-entropy above, and
\[
\int f_z^{(k)}\log f_z\,dz = -H_z\!\left(\theta^{(k)}\right) - \KL\!\left(f_z^{(k)}\,\middle\|\,
f_z\right).
\]
At $\theta = \theta^{(k)}$ its gradient is
$\int f_z\,\partial_\theta\log f_z\,dz = \partial_\theta\int f_z\,dz = 0$, so that
$\nabla Q(\theta\mid\theta^{(k)})|_{\theta=\theta^{(k)}} = \nabla\ell(\theta^{(k)})$, and a fixed point
of EM is a stationary point of $\ell$ \citep{EM}. Writing $-H_z(\theta)$ sets $\theta^{(k)} = \theta$
before the maximisation, so that the distribution of the missing sizes moves with the parameters;
its gradient $-\partial_\theta H_z(\theta)$ is not zero in general (for instance in the scale), and it
adds to the score of \eqref{eq:observed}.
\end{proof}

\begin{corollary}[The scale is shrunk]
\label{cor:shrink}
For the normal distribution with fixed location, the maximiser of \eqref{eq:entropy} in the
variance is $\hat\sigma^2_H = \frac{1}{T}\sum_{O_t=1}e_t^2 = \frac{T_1}{T}\hat\sigma^2$, where
$\hat\sigma^2 = \frac{1}{T_1}\sum_{O_t=1}e_t^2$ is the fixed point of EM.
\end{corollary}

\begin{proof}
With $H_z = \frac12\log(2\pi e\sigma^2)$, the derivative of \eqref{eq:entropy} in $\sigma^2$ is
$-\frac{T_1}{2\sigma^2} + \frac{\sum e_t^2}{2\sigma^4} - \frac{T_0}{2\sigma^2}$. EM iterates
$\sigma^2_{(k+1)} = (\sum_{O_t=1}e_t^2 + T_0\sigma^2_{(k)})/T$, whose fixed point is $\hat\sigma^2$.
\end{proof}

\begin{remark}
Two further consequences follow from the dependence of $H_z$ on the parameters. For the log-normal
iETS(M,N,N), $-H_z = -\log\hat l_{t-1} + \frac{\sigma^2}{2} - \frac12\log(2\pi e\sigma^2)$ grows
without bound as the fitted level goes to zero, which pulls the location. With a dynamic model for
the scale \citep{ADAM}, the zeros alone reward a scale going to zero (to infinity for
the Gamma distribution, whose entropy decreases in the scale for large values), and the likelihood
is unbounded. Note also that, with $H = -\E\log f$, the sign of the entropy term in
\citet[eqs.~29 and B.6]{iETS} should be negative, as in \eqref{eq:entropy}.
\end{remark}

% =============================================================================
\section{Linear Gaussian models: the exact likelihood}
\label{sec:linear}

Consider the pure additive model \eqref{eq:additive} with $\epsilon_t\sim\mathcal N(0,\sigma^2)$ in
its single-lag form, and let $\vect x_0$ be a parameter. Define $\vect a_t = \E(\vect x_{t-1}\mid
Y_{t-1})$ and $\sigma^2\vect P_t = \Var(\vect x_{t-1}\mid Y_{t-1})$, the conditional mean and
covariance of the state entering period $t$.

\begin{proposition}[The filter of the SSOE model]
\label{prop:kalman}
With $\vect a_1 = \vect x_0$ and $\vect P_1 = \vect 0$, the recursions
\begin{align}
O_t = 1:&\quad v_t = z_t - \vect w'\vect a_t,\quad f_t = \vect w'\vect P_t\vect w + 1,\quad
\vect k_t = (\vect F\vect P_t\vect w + \vect g)/f_t,\nonumber\\
&\quad \vect a_{t+1} = \vect F\vect a_t + \vect k_t v_t,\quad
\vect P_{t+1} = \vect F\vect P_t\vect F' + \vect g\vect g' - \vect k_t\vect k_t' f_t,
\label{eq:kalman}\\
O_t = 0:&\quad \vect a_{t+1} = \vect F\vect a_t,\quad
\vect P_{t+1} = \vect F\vect P_t\vect F' + \vect g\vect g',\nonumber
\end{align}
give $z_t\mid Y_{t-1}\sim\mathcal N(\vect w'\vect a_t,\sigma^2 f_t)$ at the observed periods, and the
observed-data log-likelihood of the sizes is
\begin{equation}
\ell(\theta,\sigma^2) = -\frac12\sum_{O_t=1}\left[\log(2\pi\sigma^2 f_t) +
\frac{v_t^2}{\sigma^2 f_t}\right],\qquad
\hat\sigma^2 = \frac{1}{T_1}\sum_{O_t=1}\frac{v_t^2}{f_t}.
\label{eq:kalmanlik}
\end{equation}
\end{proposition}

\begin{proof}
The model is linear and Gaussian, so $(\vect x_{t-1}, z_t, \vect x_t)$ is jointly Gaussian given
$Y_{t-1}$. With $\vect x_{t-1}\mid Y_{t-1}\sim\mathcal N(\vect a_t,\sigma^2\vect P_t)$ and $\epsilon_t$
independent of $\vect x_{t-1}$: $\E(z_t\mid Y_{t-1}) = \vect w'\vect a_t$,
$\Var(z_t\mid Y_{t-1}) = \sigma^2(\vect w'\vect P_t\vect w + 1)$, and, since
$\vect x_t = \vect F\vect x_{t-1} + \vect g\epsilon_t$ shares $\epsilon_t$ with $z_t$,
$\Cov(\vect x_t, z_t\mid Y_{t-1}) = \sigma^2(\vect F\vect P_t\vect w + \vect g)$. The conditional
distribution of $\vect x_t$ given $z_t$ gives the update at the observed periods; at the unobserved
ones $\vect x_t$ is only propagated. The decomposition of Proposition~\ref{prop:observed} gives
\eqref{eq:kalmanlik}, and $\sigma^2$ concentrates out because $f_t$ and $\vect k_t$ do not depend on
it.
\end{proof}

The filter is the Kalman filter for a state space model with correlated disturbances
\citep{Harvey1989,DurbinKoopman2012}, written for the single source of error.

\begin{corollary}[Complete data]
\label{cor:complete}
If $\vect P_t = \vect 0$ and $O_t = 1$, then $\vect k_t = \vect g$, $f_t = 1$ and
$\vect P_{t+1} = \vect 0$. Hence, when all periods are observed, \eqref{eq:kalman} is the SSOE
recursion and \eqref{eq:kalmanlik} the likelihood of the one-step errors.
\end{corollary}

\begin{proof}
$\vect k_t = \vect g/1$ and $\vect P_{t+1} = \vect g\vect g' - \vect g\vect g' = \vect 0$.
\end{proof}

Corollary~\ref{cor:complete} is the formal statement of the property that makes the usual
estimation of ETS exact \citep{OrdKoehlerSnyder1997}. With unobserved periods it fails: after
$j-1$ unobserved periods following a period with $\vect P = \vect 0$,
\begin{equation}
\vect P = \sum_{i=1}^{j-1}\vect F^{i-1}\vect g\vect g'\vect F'^{\,i-1}, \qquad
f = 1 + \sum_{i=1}^{j-1}c_i^2 \equiv \kappa_j, \qquad c_i = \vect w'\vect F^{i-1}\vect g,
\label{eq:kappa}
\end{equation}
and the update at the next observed period leaves $\vect P_{t+1}\neq\vect 0$ unless the observation
reveals the whole drift. The state uncertainty then persists and decays over the following
observations, and the gain differs from $\vect g$ as long as it does.

\begin{proposition}[The current estimator is not the maximum likelihood estimator]
\label{prop:notmle}
With unobserved periods and $\vect g\neq\vect 0$, the recursion \eqref{eq:recursion} with the
one-step density differs from \eqref{eq:kalmanlik} in three respects: the variance of the error
at an observed period is $\sigma^2 f_t\geq\sigma^2$ rather than $\sigma^2$; the gain of the update is
$\vect k_t$ rather than $\vect g$, so that the fitted states and the errors differ; and these
differences persist after the unobserved periods through $\vect P_t$. Its maximiser is therefore not
the maximum likelihood estimator.
\end{proposition}

\begin{proof}
By Proposition~\ref{prop:kalman} the predictive distribution is $\mathcal N(\vect w'\vect a_t,
\sigma^2 f_t)$, and $f_t > 1$ whenever $\vect P_t\neq\vect 0$, which holds at the first observed period
after an unobserved one with $c_1\neq 0$, and at the following ones until $\vect P_t$ returns to zero.
The conditional mean $\vect a_t$ follows the gain $\vect k_t\neq\vect g$ at those periods, while
\eqref{eq:recursion} uses $\vect g$.
\end{proof}

\begin{example}
For the local level model, ETS(A,N,N), with $\vect P = 0$ before an interval of $j$ periods, the
gain at its end is
\begin{equation}
\alpha_j = \frac{(j-1)\alpha^2 + \alpha}{(j-1)\alpha^2 + 1},
\label{eq:gain}
\end{equation}
which is $\alpha$ for $j = 1$ and tends to one as $j$ grows: after a long interval the level should
jump almost to the new observation, while \eqref{eq:recursion} moves it by $\alpha$ times the error.
The remaining variance $P_{t+1} = (j-1)\alpha^2 + \alpha^2 - \alpha_j^2\kappa_j$ is positive for
$j > 1$ and $0<\alpha<1$, and each following observed period multiplies it by
$(1-\alpha)^2/(1+P_t)$.
\end{example}

\begin{remark}
Two approximations of \eqref{eq:kalman} are of interest because they keep the structure of the
SSOE recursion. Resetting $\vect P$ to zero after each observed period gives the variance
$\kappa_{j_t}$ of \eqref{eq:kappa} at the observed periods and, with the gain \eqref{eq:gain} (or
its multivariate counterpart), the first-order correction of the states. Keeping only $\kappa_{j_t}$
and the gain $\vect g$ corrects the likelihood but not the states, and underestimates the variance
at the observed periods that follow a long interval.
% TODO: illustration of the size of the three versions on simulated series.
\end{remark}

% =============================================================================
\section{Pure multiplicative ETS of ADAM}
\label{sec:multiplicative}

\begin{lemma}[Linearity in logarithms]
\label{lem:loglinear}
In the pure multiplicative ETS of ADAM, with $u_t = \log(1+\epsilon_t)$, the logarithms of the
components satisfy
\begin{equation}
\log z_t = \vect w'\log\vect v_{t-\vect l} + u_t, \qquad
\log\vect v_t = \vect F\log\vect v_{t-\vect l} + \vect g\,u_t ,
\label{eq:loglinear}
\end{equation}
with the $\vect w$, $\vect F$ and $\vect g$ of the corresponding additive model. In the conventional
form, the innovation of component $k$ is $L_k(\epsilon_t) = \log(1+g_k\epsilon_t)$ instead of
$g_k u_t$, and the model is linear in logarithms with a non-Gaussian innovation vector.
\end{lemma}

\begin{proof}
Take logarithms of the measurement $z_t = l_{t-1}b_{t-1}^{\phi}s_{t-m}(1+\epsilon_t)$ and of the
updates, for instance $\log l_t = \log l_{t-1} + \phi\log b_{t-1} + \alpha u_t$ and
$\log b_t = \phi\log b_{t-1} + \beta u_t$ in the ADAM form, and
$\log l_t = \log l_{t-1} + \phi\log b_{t-1} + \log(1+\alpha\epsilon_t)$ in the conventional one.
\end{proof}

With the log-normal distribution, $u_t\sim\mathcal N(-s/2, s)$ so that $\E(1+\epsilon_t) = 1$, and
\eqref{eq:loglinear} is a linear Gaussian model for $\log z_t$ with a known mean of the errors.

\begin{proposition}[Exact likelihood of the log-normal ADAM ETS]
\label{prop:lognormal}
For the pure multiplicative ETS of ADAM with the log-normal distribution, the filter
\eqref{eq:kalman} applied to $\log z_t$, with $u_t = -s/2 + \eta_t$, $\eta_t\sim\mathcal N(0,s)$,
the mean $-s\vect g/2$ added to the propagation of $\vect a_t$ and $-s/2$ to the prediction of
$\log z_t$, gives the exact observed-data likelihood
\[
\ell(\theta, s) = -\frac12\sum_{O_t=1}\left[\log(2\pi s f_t) + \frac{v_t(s)^2}{s f_t}\right]
- \sum_{O_t=1}\log z_t ,
\]
where $v_t(s) = v_t^0 + s\,b_t$ is linear in $s$ because $f_t$ and the gains do not depend on $s$.
The maximum likelihood estimate of the scale is the positive root of
\begin{equation}
\left(\sum_{O_t=1}\frac{b_t^2}{f_t}\right)s^2 + T_1 s - \sum_{O_t=1}\frac{(v_t^0)^2}{f_t} = 0 .
\label{eq:lognormalscale}
\end{equation}
\end{proposition}

\begin{proof}
The filter is that of Proposition~\ref{prop:kalman} for the linear Gaussian model
\eqref{eq:loglinear}, with the Jacobian of the transformation. The predictions $\vect w'\vect a_t$
are linear in the known means $-s/2$ with coefficients that depend on the gains only, which gives
$v_t(s) = v_t^0 + s b_t$. Setting the derivative in $s$ to zero and multiplying by $2s^2$ gives
\eqref{eq:lognormalscale}.
\end{proof}

When $\vect P$ is zero at the previous observed period, $f_t = \kappa_{j_t}$ and
$b_t = a_{j_t}/2$ with $a_j = 1 + \sum_{i<j}c_i$, so that
$\log(z_t/\hat\mu_t)\sim\mathcal N(-s a_{j_t}/2, s\kappa_{j_t})$ and
\begin{equation}
\hat s = \frac{2\left(\sqrt{T_1^2 + \sum\frac{a_{j_t}^2}{\kappa_{j_t}}\sum\frac{u_t^2}{\kappa_{j_t}}}
- T_1\right)}{\sum a_{j_t}^2/\kappa_{j_t}} .
\label{eq:lognormalscalereset}
\end{equation}

\begin{corollary}[The scale of the log-normal distribution]
\label{cor:lognormalscale}
Without unobserved periods the maximum likelihood estimate of $s$ is
$\hat s = 2(\sqrt{1+m}-1)$ with $m = \frac{1}{T_1}\sum u_t^2$. The estimator
$2(1-\sqrt{1-m})$, the root of $s^2 - 4s + 4m = 0$, is not the maximum likelihood one; it
overestimates $s$ (by about 35\% at $s = 0.5$) and is undefined for $m > 1$.
\end{corollary}

\begin{proof}
Set $f_t = 1$ and $b_t = 1/2$ in \eqref{eq:lognormalscale}: $\frac{T_1}{4}s^2 + T_1 s - T_1 m = 0$,
that is $s^2 + 4s - 4m = 0$. For $s = 0.5$, $m = s + s^2/4 = 0.5625$ and $2(1-\sqrt{1-m}) = 0.677$.
\end{proof}

% =============================================================================
\section{Other pure multiplicative models: an approximate predictive density}
\label{sec:approximate}

For the Gamma and the inverse Gaussian distributions, and for the conventional ETS with any of the
three, Lemma~\ref{lem:loglinear} still gives a linear model for the logarithms of the states, but
with non-Gaussian innovations $L_k(\epsilon_t)$, and the measurement $z_t = \mu_t(1+\epsilon_t)$
keeps the one-step distribution of $1+\epsilon_t$. We approximate the conditional distribution of
the logarithm of the location by a normal one and keep the one-step density exact.

\subsection{The drift of the states}

Let the states at the previous observed period be known (the reset approximation of
Section~\ref{sec:linear}). Over $j_t - 1$ unobserved periods the logarithm of the location drifts by
\begin{equation}
d_t = \log\frac{\mu_t}{\hat\mu_t} = \sum_{i=1}^{j_t-1}\sum_k c_{i,k}\,L_k(\epsilon_{t-i}),
\qquad c_{i,k} = \vect w'\vect F^{i-1}\vect e_k ,
\label{eq:drift}
\end{equation}
where $\vect e_k$ is the $k$-th unit vector, so that $c_{i,k}$ are the coefficients $c_i$ of
\eqref{eq:kappa} for the persistence of the component $k$ alone. Its mean and variance are
\begin{equation}
m_t = \sum_{i<j_t}\vect c_i'\boldsymbol\mu_L, \qquad
w_t = \sum_{i<j_t}\vect c_i'\boldsymbol\Sigma_L\vect c_i,
\label{eq:driftmoments}
\end{equation}
with $\boldsymbol\mu_L$ and $\boldsymbol\Sigma_L$ the mean and the covariance of the innovation
vector. These depend on the distribution and the scale. We compute them by Gauss--Legendre
quadrature on the quantile function $Q$ of $1+\epsilon$:
$\E h(\epsilon) = \int_0^1 h(Q(u)-1)\,du \approx \sum_q \omega_q h(Q(u_q)-1)$
\citep{GolubWelsch1969}. The integrands $\log(1-g+gQ(u))$ and $g\log Q(u)$ have at most logarithmic
singularities at the ends of the interval, and 32 nodes are ample for the accuracy needed here.

\subsection{The mixture density}

With $d_t\sim\mathcal N(m_t, w_t)$, the predictive density of the size is
\begin{equation}
\tilde f(z_t) = \int f_z\!\left(z_t\mid\hat\mu_t e^{d}\right)\phi(d; m_t, w_t)\,dd
\approx \sum_{q=1}^{Q}\frac{\omega_q}{\sqrt\pi}\,
f_z\!\left(z_t\mid\hat\mu_t e^{m_t + \sqrt{2w_t}x_q}\right),
\label{eq:mixture}
\end{equation}
with the Gauss--Hermite nodes $x_q$ and weights $\omega_q$ \citep{LiuPierce1994}. It is exact for
$j_t = 1$. For the log-normal distribution the mixture is log-normal again,
$\log(z_t/\hat\mu_t)\sim\mathcal N(-s/2 + m_t, s + w_t)$, and no quadrature is needed.

The normal approximation of the drift is the only approximation in \eqref{eq:mixture}. Against the
exact density of the size after an interval, computed by numerical convolution, the largest
Kullback--Leibler divergence per observation over a grid of $s\in[0.05,1]$, $\alpha\in[0.1,0.6]$
and $j\in[2,20]$ for iETS(M,N,N) was 0.0007 for the Gamma, 0.0033 for the log-normal and 0.0038 for
the inverse Gaussian distribution, against 2.36, 5.05 and 31.5 for the one-step density.
% TODO: describe the computation of the exact density and report the grid in a table.

\subsection{The scale}

ADAM estimates the scale of the Gamma and the inverse Gaussian distributions by moments rather than
by maximum likelihood \citep{ADAM}: $\hat s = \frac{1}{T_1}\sum e_t^2$ for the Gamma
and $\hat s = \frac{1}{T_1}\sum e_t^2/(1+e_t)$ for the inverse Gaussian, with the relative errors
$e_t = z_t/\hat\mu_t - 1$ and $\Var(1+\epsilon_t) = s$. After an interval,
$1 + e_t = (1+\epsilon_t)D_t$ with $D_t = e^{d_t}$ independent of $\epsilon_t$.

\begin{proposition}[Moment equations with the drift]
\label{prop:moments}
With $\E(1+\epsilon_t) = 1$ and $\Var(1+\epsilon_t) = s$,
\begin{align*}
\E(e_t^2) &= (1+s)\,\E(D_t^2) - 2\,\E(D_t) + 1,\\
\E\!\left(\frac{e_t^2}{1+e_t}\right) &= \E(D_t) - 2 + (1+s)\,\E(D_t^{-1})
\quad\text{(inverse Gaussian)},
\end{align*}
the latter because $\E[(1+\epsilon)^{-1}] = 1 + s$ for the inverse Gaussian distribution with mean
one and dispersion $s$. With the normal drift, $\E D_t = e^{m_t + w_t/2}$, $\E D_t^2 = e^{2m_t+2w_t}$
and $\E D_t^{-1} = e^{-m_t + w_t/2}$.
\end{proposition}

\begin{proof}
Expand $e_t^2 = (1+e_t)^2 - 2(1+e_t) + 1$ and $e_t^2/(1+e_t) = (1+e_t) - 2 + (1+e_t)^{-1}$, use the
independence of $\epsilon_t$ and $D_t$, and the moments of the inverse Gaussian and the log-normal
distributions.
\end{proof}

Because $m_t$ and $w_t$ depend on $s$ through $\boldsymbol\mu_L$ and $\boldsymbol\Sigma_L$, the
estimating equations obtained by replacing the expectations by sample means,
\begin{equation}
s = h(s) = \frac{\sum e_t^2 - \sum\left(\E D_t^2 - 2\E D_t + 1\right)}{\sum\E D_t^2}
\quad\text{(Gamma)},
\label{eq:fixedpoint}
\end{equation}
and the analogous one for the inverse Gaussian, are solved by the fixed-point iteration
$s^{(n+1)} = h(s^{(n)})$ from the one-step estimate $s^{(0)}$. Without intervals $h$ does not depend
on $s$ and the iteration returns the one-step estimate at once. The dependence of $h$ on $s$ is weak,
of the order of $(j_t - 1)\alpha^2 s$ in iETS(M,N,N), so that $h$ is a contraction for the values of
the parameters of interest and a few iterations suffice. For the log-normal distribution the
estimate is \eqref{eq:lognormalscalereset} with $a_t = 1 - 2m_t/s$ and $\kappa_t = 1 + w_t/s$,
iterated in the same way.

\begin{remark}
On simulated errors with the drift of iETS(M,N,N) these equations recover the scale without bias.
Within a fitted model the estimate is still biased upwards, because the states after an interval
are updated with the gain $\vect g$ rather than with the conditional mean, so that part of the drift
leaks into the following errors (Proposition~\ref{prop:notmle}). The approximate filter of the next
remark addresses both.
\end{remark}

\begin{remark}[An approximate filter]
The reset approximation can be lifted as in Section~\ref{sec:linear}: carry a covariance
$\vect P_t$ of the logarithms of the states, propagate it over the unobserved periods with
$\boldsymbol\Sigma_L$, use $d_t\sim\mathcal N(m_t, \vect w'\vect P_t\vect w)$ in \eqref{eq:mixture},
and update the states with the linear projection of the logarithm of the state on $\log z_t$, which
reduces to the conventional update when $\vect P_t = \vect 0$. This is a moment-matching filter in the
spirit of the extended and unscented filters \citep{DurbinKoopman2012}.
% TODO: write the update explicitly and check that it reduces to the SSOE update at P = 0.
\end{remark}

% =============================================================================
\section{Prediction intervals}
\label{sec:intervals}

A forecast $h$ steps ahead from the end of the sample is the prediction after an unobserved stretch
of $h-1$ periods, so the results above give the predictive distributions. For the linear Gaussian
models the variance is $\sigma^2(\vect w'\vect F^{h-1}\vect P_{T+1}\vect F'^{\,h-1}\vect w +
\kappa_h)$: the state covariance at the end of the sample, zero for complete data, adds to the usual
multistep variance when the last periods were unobserved. For the pure multiplicative ADAM ETS with
the log-normal distribution, Proposition~\ref{prop:lognormal} gives the exact interval
$z_{T+h}/\hat\mu_{T+h}\sim\log\mathcal N(-s a_h/2, s\kappa_h)$ around the point forecast with zero
errors, when $\vect P_{T+1} = \vect 0$.

For a model with an occurrence part, an interval of the demand with coverage $c$ that counts the
zeros as covered is an interval of the sizes with coverage $(c - (1-p_t))/p_t$, so its bounds are
quantiles of the sizes around their forecast $\hat\mu_t$. Centring them on the forecast of the
demand, $p_t\hat\mu_t$, which is the conditional mean of $y_t$ but not a location of the sizes,
places the upper bound well below the corresponding quantile of the demand.

% =============================================================================
\section{Computation and implementation}
\label{sec:computation}

\paragraph{Cost.} The filter \eqref{eq:kalman} costs $O(n^2)$ operations per period for $n$
states, against $O(n)$ for the SSOE recursion, but only while $\vect P_t\neq\vect 0$: it can switch
to the SSOE recursion whenever the covariance has decayed below a tolerance. The reset
approximation needs only the coefficients $c_i$ up to the longest interval, which are those of the
analytical multistep variance of ADAM \citep{ADAM}, computed once per evaluation of
the likelihood. The quadrature of Section~\ref{sec:approximate} adds 32 evaluations of the quantile
function per iteration of the scale and per evaluation, and the mixture 16 evaluations of the
one-step density per observation that follows an interval.

\paragraph{Reproducibility.} The quadrature rules are fixed sets of numbers, so that two
implementations evaluate the same likelihood. The coefficients $c_i$ are computed by one routine in
C++ shared by the R and the Python versions of the package, and the sums are accumulated in the
same order in both.

\paragraph{Implementation.} The \texttt{smooth} package (R and Python) implements, on its
development branch: the observed-data likelihood without the entropy (Proposition~\ref{prop:observed});
the reset approximation with the variance $\kappa_{j_t}$ for the pure additive models with the
normal distribution and the exact gap density of the pure multiplicative ADAM ETS with the
log-normal one, with the scale \eqref{eq:lognormalscalereset}; the mixture \eqref{eq:mixture} with
the moment equations of Proposition~\ref{prop:moments} for the other pure multiplicative models; the
corrected scale of the log-normal distribution (Corollary~\ref{cor:lognormalscale}); and the
prediction intervals of Section~\ref{sec:intervals}. The filter \eqref{eq:kalman}, with the gain it
implies, is not implemented: it changes the fitted states, and with them the fitter of the package.
% TODO: decide whether the paper presents the filter as the method and the reset as its
% approximation, and implement the filter before the evaluation.

% =============================================================================
\section{Conclusions}
\label{sec:conclusions}

The estimation of exponential smoothing by its own recursion and the likelihood of the one-step
errors is exact only because, with complete data, the states are known functions of the data. When
some periods are unobserved, as with the zeros of intermittent demand or with missing values, the
states become uncertain, and the estimator of ADAM and iETS is no longer the maximum likelihood one.
Adding the differential entropy of the unobserved sizes does not repair this, and biases the scale.
For the linear Gaussian models and for the pure multiplicative ETS of ADAM with the log-normal
distribution, the exact likelihood is given by a filter whose gain and variance vary after the
unobserved periods and reduce to those of the SSOE recursion with complete data. For the other pure
multiplicative models a Gaussian approximation of the drift of the logarithms of the states gives an
accurate predictive density. The mixed models, with additive and multiplicative components, are
linear in neither space and are left for further work, as are the empirical evaluation of the
estimators and the forecasting performance of the corrected models.
% TODO: evaluation on simulated and real intermittent series.

\begin{thebibliography}{99}

\bibitem[Dempster et~al.(1977)]{EM}
Dempster, A.~P., Laird, N.~M. and Rubin, D.~B. (1977). Maximum likelihood from incomplete data via
the EM algorithm. \emph{Journal of the Royal Statistical Society, Series B}, 39(1), 1--38.

\bibitem[Durbin and Koopman(2012)]{DurbinKoopman2012}
Durbin, J. and Koopman, S.~J. (2012). \emph{Time Series Analysis by State Space Methods}, 2nd
edition. Oxford University Press.

\bibitem[Golub and Welsch(1969)]{GolubWelsch1969}
Golub, G.~H. and Welsch, J.~H. (1969). Calculation of Gauss quadrature rules. \emph{Mathematics of
Computation}, 23(106), 221--230.

\bibitem[Harvey(1989)]{Harvey1989}
Harvey, A.~C. (1989). \emph{Forecasting, Structural Time Series Models and the Kalman Filter}.
Cambridge University Press.

\bibitem[Hyndman et~al.(2008)]{Hyndman2008}
Hyndman, R.~J., Koehler, A.~B., Ord, J.~K. and Snyder, R.~D. (2008). \emph{Forecasting with
Exponential Smoothing: The State Space Approach}. Springer.

\bibitem[Liu and Pierce(1994)]{LiuPierce1994}
Liu, Q. and Pierce, D.~A. (1994). A note on Gauss--Hermite quadrature. \emph{Biometrika}, 81(3),
624--629.

\bibitem[Ord et~al.(1997)]{OrdKoehlerSnyder1997}
Ord, J.~K., Koehler, A.~B. and Snyder, R.~D. (1997). Estimation and prediction for a class of
dynamic nonlinear statistical models. \emph{Journal of the American Statistical Association},
92(440), 1621--1629.

\bibitem[Snyder(1985)]{Snyder1985}
Snyder, R.~D. (1985). Recursive estimation of dynamic linear models. \emph{Journal of the Royal
Statistical Society, Series B}, 47(2), 272--276.

\bibitem[Svetunkov(2023)]{ADAM}
Svetunkov, I. (2023). \emph{Forecasting and Analytics with the Augmented Dynamic Adaptive Model
(ADAM)}. Chapman and Hall/CRC.

\bibitem[Svetunkov and Boylan(2023)]{iETS}
Svetunkov, I. and Boylan, J.~E. (2023). iETS: State space model for intermittent demand
forecasting. \emph{International Journal of Production Economics}, 265, 109013.

\end{thebibliography}

\end{document}
